% !TeX root = ../Western Philosophy.tex
\section{霍布斯哲学}

\subsection{生平简介（1588）}

在年轻的时候几乎没有与哲学的交集，开始进行人文写作的标志是1640年《法的要素》。对Discart进行了批评，关于物理世界的解释是一样的，但是关于其对心灵世界的解释是及其不认同的。

1651年发表了最著名的著作《利维坦》，当时流亡法国，在伦敦出版。

\subsection{唯物的经验论}

霍布斯提出唯物经验论，提出了感觉的基础地位，同时认为世界上不存在认知或者感知中的实体。有以下观点：

\begin{itemize}
    \item 感觉是被动接受的过程
    \item 感觉是一切观念的基础：记忆，想象和推理等都基于感觉。
    \item 两种知识：事实知识和后果知识
\end{itemize}

\subsection{唯名论与语言}

首先是对于唯名论和语言，提出了对共相（universal）问题的处理。

\begin{fancybox}{共相}
    共相可以理解为一种公共的性质。譬如一间教室里面有很多白色衣服的同学，共相实在论认为为存在一个“白色性”的实体，每一个白色衣服都是实体的“例示”。对于唯名论者，如何在不承认一个实体存在的情况下保留事物的共有属性。
\end{fancybox}

唯名论相对于实在论，反对超出具体事物的分有属性。霍布斯对于共相问题的处理如下：

\begin{itemize}
    \item 不存在共相
    \item 所有词语都是名字，命名一个或者多个个体
    \item 句子的真值在于名字命名关系的恰当重叠（eg：人是生物，所有人都是正义的）
\end{itemize}

霍布斯对于语言哲学进行了一般论述（可以看到这个系统化程度和体系化程度低于欧陆）：

\begin{itemize}
    \item 语言是人和动物的最大区别（区别于Discart的心灵能力）
    \item 语言的功能（虽然是不知所云的分类）：记录知识，咨询与讲授，声明意愿，娱乐。
    \item 语言是筹码，而非实值：Words are wise men’s counters，they do bur reckon by them;but they are the money of fools. 
\end{itemize}

\subsection{心灵与意志}

Hobbes的心灵哲学首先在于对于Dsicart实体二元论的批判，他对心身交互问题的批评在于：无身体的实体如同圆的方。

关于自由意志的兼容论（英国哲学的经验论传统和兼容论传统）：

\begin{itemize}
    \item 意志知识出现在思虑最后一步的欲望，人和动物没有本质差异。
    \item 人和动物的欲望都可以机械解释（唯物论图景下的心灵状态）
    \item 自由与必然可以兼容（出自意志，所以自由；出自原因，所以必然）
    
    如果意志则自由行动（free to do it if he will）vs 自由地意志（free to will），
\end{itemize}

\begin{fancybox}{一些实践哲学名词}
    deliberation：对于不同因素的权衡
    
    reason：一种faculty，用于做出理性决策（很多人认为这是人和动物的区别）

    decision：权衡和决定的结果
\end{fancybox}

\subsection{政治哲学}

政治哲学探讨几个问题：Collective question 

\subsubsection{自然状态}

三大原因导致纷争：竞争，不自信，荣誉（这个分类没什么逻辑，就是这样分的）

由此形成著名的\textbf{所有人对所有人的战争（a war of every man against every man）}人生在自然状态下是各种不堪的（solitary，poor，nasty，brutish and short）

\subsubsection{自然法要求创造主权者}

所谓自然状态指的是“理性自利”。根据理性计算保持自身发展

\begin{itemize}
    \item 自然法vs自然权利（Natural right）
    \item （理性自利的自然法要求）所有人让步，把除了自我防卫之外的权利让渡给一个依靠惩罚执法的中央权力。
    \item 自然法也直接要求人们遵守各种契约，这也需要强权，保证违约成本高于违约收益。Covnents without the sword，are but words，cannot secure a man at all。
\end{itemize}

\subsubsection{主权者的构成与权限}

利维坦中的主权者，主权者可以是个人（如独裁者）也可以是群体（如议会）

主权者拥有广泛的决定权，可以决定内外法治，决定官员的任命和赏罚，可以决定观点和教条被允许，在宗教上是至高的。

在法律不规定的方面（具体的生活方式），臣民有自由。

主权者无权要求臣民进行自我伤害或者毁灭（即使被定罪，这样的情况臣民也拥有反击的自由）。

\begin{fancybox}{Hobbes利维坦的问题}
    \begin{itemize}
        \item Hobbes的政治哲学立论需要做哪些立论预设？（权力包含什么内容？生命权是可以让渡的吗？）
        \item 如果把Hobbes的政治哲学付诸实践，可能会遇到什么问题？（如果有超级英雄/能力强于其他个体的人，自然状态下权利大于在利维坦中的权利，则这种人在Hobbes哲学下不需要利维坦；如果有多个利维坦，利维坦具有人格吗，有无自然状态）
        \item 诉诸利维坦（或者权力体）是走出自然状态的唯一解吗？（亚当斯密，休谟主张道德和共情来走出自然状态）
    \end{itemize}
\end{fancybox}
